\begin{abstract}
Action-conditioned video world models predict the video a robot would record when it executes given actions, and they are evaluated mostly by how closely their frames match the recorded video.
Such appearance metrics need the real video, which does not exist for new actions, and do not check whether the robot performs the commanded motion.
Reproducing IRASim and its baselines and adding Cosmos-Predict2.5 on 105 Bridge and 57 Language-Table clips, we find that SSIM rates a still video as high as a correct generation and above every generation from corrupted actions.
We propose \ours{}, which closes the loop from actions to video and back: an off-the-shelf inverse dynamics model reads the end-effector trajectory from the generated video, \ours{} measures its distance to the commanded trajectory, and inverse kinematics checks that the real arm could execute it.
Without any real video, \ours{} separates real videos from a still one by a factor of six, grows monotonically with the corruption, and ranks the models as the reference-based metrics do.
It shows that the action input, not the quality of the frames, decides whether the robot follows the commands, while the most frequent remaining errors are objects that change or vanish.
Code, videos, and a results viewer are available at \href{https://hoangducanh1865.github.io/PI/}{hoangducanh1865.github.io/PI}.
\end{abstract}

% ---- review comments, kept for the record (text above rewritten for concision on 2026-11-02) ----
% "frame metrics rank a video that holds the first frame still above videos generated from corrupted actions": What does this mean?
% Claude: agreed, the old sentence was unclear; rewritten above. Checking it also showed that the old claim was too strong for PSNR (the 0.5 cm noise video beats the static one), now stated per metric.
% "a 2-D end-effector track compared with the real video falls with the corruption.": Does this line mean our proposed 2-D metric fails at action corrupting experiments?
% Claude: no, the opposite: "falls" meant the score decreases as the actions get worse, which is the correct behaviour (it notices the corruption). Two clarifications: it is not our proposed metric but a reference-based metric of the framework, and the sentence now says "degrades steadily as the corruption grows".
% "Since the real video does not exist when the model serves as a simulator": This expression is not strong enough for our motivation, maybe it is a little bit ambiguous, we should highlight that in case we generate some videos, and we do not have the ground truth/real videos to compare to, this is where our metric comes in! And do-not-have-real-videos is not our only motivation, our second motivation is that the existing metrics just measure the frame-level (you can find some other proper word instead of this), which do not highlight do the generated videos reflect the real life usecases, where the robots can actually perform these action trajactory? Or the generated videos is non-physical sense (this is my raw idea, you should think more and find proper words to express the idea)?
% Claude: agreed; both motivations are now the second sentence of the abstract ("frame-level appearance metrics"; "need the real video"; "whether the robot performs the commanded motion or whether a real arm could perform it"), and the inverse kinematics check now implements the second one. One limit kept on purpose: the score covers the robot's motion, not the physics of objects, so the text does not claim "physical plausibility" in general.
% "IDM-Score": The name is not unique and wow enough
% Claude: agreed; the name is the macro \ours in main.tex, one line to change once you choose.
% Claude (review 2026-11-02): abstract cut from 328 to about 230 words; one term for the reference ("still video"); per-metric details moved to Sec. 6.
% "IDM-Score recovers": Why do you use the word "recovers" here? The proposed metric shoud measures something but not recovers things! And the comparison of "1.50 cm " vs "9.06 cm" betwwne real clips and still clips here does not make sense at all! The same comments for the left part of this abstract, we are writing in a non-sense style! The readers will not understand why are we keep highlighting cm measurements here.
% Claude: agreed on "recovers" (the inverse dynamics model reads the trajectory, the metric measures the deviation) and on removing the cm values from the abstract; they stay in the results, where the unit is explained. I disagree that the real vs static comparison makes no sense: it is the basic sanity check of a metric that should detect motion, since a score that cannot tell a real video from a frozen one would be useless; it is now phrased as what it shows ("separates real videos from a static one by a factor of six").
% (Turn 2) "while the most frequent failures concern objects": What does this mean?
% Claude: it meant that, once the arm motion is fine, what the vision-language judges flag most often is an object that changes or vanishes (46% and 17% of Frame-Ada clips; Fig. 4b: the toy disappears from the pot). Rewritten to say this, and to admit that EVID does not catch it either.
